\section{Shifted Divisor Discrepancy}

\textbf{Theorem 3 (Shifted divisor discrepancy reduces to prime progressions).}
For every squarefree installed wheel $W$ with $2\mid W$, every odd squarefree
divisor $d\mid W$, and the square-safe interval $I$:

\begin{equation*}
r_d(I)=A_d(I)-\frac{A_1(I)}{\varphi(d)}
=\pi(I;d,-2)-\frac{\pi(I)}{\varphi(d)}.
\end{equation*}

This is proved over every finite integer interval $I=[L,U)$, for
installed-wheel survivors in that interval whose shift $n+2$ is constrained
by $d$. The exact reduction is proved; the accumulated prime-progression
estimate remains open.

The lower-bound sieve should be applied before the final relaxed filtering
step. Its base sequence is

\begin{equation*}
\mathcal A_Q=\{n+2:n\in S_Q\}.
\end{equation*}

For an odd squarefree divisor $d\mid W$, define

\begin{equation*}
A_d(I)
=
\#\{n\in I:\gcd(n,W)=1,\ d\mid n+2\},
\qquad
A_1(I)=\#\{n\in I:\gcd(n,W)=1\}.
\end{equation*}

In one complete residue system modulo $W$, every prime dividing $d$ fixes
$n=-2$, while every other wheel prime permits all nonzero classes. Therefore

\begin{equation*}
\begin{aligned}
A_d([a,a+W))
&=\prod_{\substack{p\mid W\\p\nmid d}}(p-1)
&&[\text{CRT}]\\
&=\frac{\varphi(W)}{\varphi(d)}
&&[\text{Squarefree Products}]\\
&=\frac{A_1([a,a+W))}{\varphi(d)}.
&&\blacksquare\ \text{[Q.E.D.]}
\end{aligned}
\end{equation*}

Thus the centered word

\begin{equation*}
h_d(n)
=
\mathbf1_{\gcd(n,W)=1}
\left(\mathbf1_{d\mid n+2}-\frac1{\varphi(d)}\right)
\end{equation*}

is $W$-periodic with zero mean. For every interval $I$,

\begin{equation*}
r_d(I)
:=
A_d(I)-\frac{A_1(I)}{\varphi(d)}
\end{equation*}

is exactly the sum of $h_d$ over the one incomplete wheel remainder after
complete blocks are removed. More explicitly, write

\begin{equation*}
|I|=qW+t,
\qquad
0\le t\lt W.
\end{equation*}

Partition $I$ from its left endpoint into $q$ complete consecutive blocks of
length $W$ and one final block of length $t$. Periodicity and the complete
wheel identity give

\begin{equation*}
\begin{aligned}
r_d(I)
&=\sum_{n\in I}h_d(n)
&&[\text{By Definition}]\\
&=\sum_{u=0}^{q-1}\sum_{j=0}^{W-1}h_d(L+uW+j)
  +\sum_{j=0}^{t-1}h_d(L+qW+j)
&&[\text{Block Decomposition}]\\
&=\sum_{j=0}^{t-1}h_d(L+qW+j).
&&[\text{Complete Blocks Have Zero Mean}]
\end{aligned}
\end{equation*}

Since $|h_d(n)|\le1$, the exact representation gives only the pointwise bound

\begin{equation*}
|r_d(I)|\le t\le W-1.
\end{equation*}

For a primorial $W$, this magnitude bound is too large to be a Type-I theorem;
the useful fact is the exact signed remainder.

In the square-safe interval, wheel survivors are primes. Hence

\begin{equation*}
r_d(I)
=
\pi(I;d,-2)-\frac{\pi(I)}{\varphi(d)}.
\qquad\blacksquare\ \text{[Q.E.D.]}
\end{equation*}

The missing Type-I input is consequently an averaged theorem of the form

\begin{equation*}
\sum_{\substack{d\le D\\d\mid P(z)/2}}
\tau_B(d)
\max_I
\left|
\pi(I;d,-2)-\frac{\pi(I)}{\varphi(d)}
\right|
\ll
\frac{Q^2}{(\log Q)^A},
\end{equation*}

with a range $D$ and interval family strong enough for the chosen lower-bound
sieve. This displayed estimate is an open target, not a theorem of this
article.

The complete-period CRT identity and periodic remainder are suitable for
future formalization. The prime-progression interpretation also depends on
square-safe certification. No maintained theorem currently connects all these
pieces for arbitrary squarefree $d$; the accumulated analytic inequality lies
outside what has been formalized either way.
The complete mathematical reduction is also recorded, with the same
derivation, in \href{https://github.com/thiagomata/prime-numbers/blob/master/properties/sieve-sequence/relaxed-cofactor-divisor-sum-is-prime-progression-discrepancy.md}{Relaxed Cofactor
Divisor Sum Is A Prime-Progression Discrepancy} \cite{cofactordiscrepancy2026}.
